\uuid{w0NX}
\titre{Équation différentielle à second membre distributionnel}

\niveau{M1}
\module{Analyse}
\chapitre{Distributions}
\sousChapitre{Équations différentielles dans $\mathcal{D}'_+$}
\theme{Convolution avec des dérivées de Dirac}
\auteur{Q. Liard}
\datecreate{ 2026-09-03}
\organisation{CentraleSupélec}
\difficulte{4}

\contenu{

\texte{
On cherche une distribution $T$ à support dans $\mathbb{R}_+$ vérifiant
$$
T''-T=\delta'-\delta''.
$$
On utilise la convention de Fourier
$$
\widehat{u}(\xi)=\int_{\mathbb{R}}u(t)e^{-i\xi t}\,dt.
$$
Lorsque la formule d'inversion est applicable, elle s'écrit
$$
u(t)=\frac{1}{2\pi}\int_{\mathbb{R}}\widehat{u}(\xi)e^{i\xi t}\,d\xi.
$$
}

\begin{enumerate}

\item \question{Résoudre le problème de Cauchy
$$
\begin{cases}
z''-z=0,\\
z(0)=0,\\
z'(0)=1,
\end{cases}
$$
puis en déduire une distribution $G$ vérifiant $G''-G=\delta$.}

\reponse{
Les solutions de $z''-z=0$ sont de la forme
$$
z(t)=Ae^t+Be^{-t}.
$$
Les conditions $z(0)=0$ et $z'(0)=1$ donnent
$$
A+B=0,
\qquad
A-B=1.
$$
Ainsi $A=1/2$ et $B=-1/2$, donc
$$
\boxed{z(t)=\sinh t}.
$$
Comme $z(0)=0$ et $z'(0)=1$,
$$
(H z)''-H z
=H(z''-z)+\delta
=\delta.
$$
On peut donc prendre
$$
\boxed{G(t)=H(t)\sinh t}.
$$
}

\item \question{En déduire une solution de
$$
T''-T=\delta'-\delta''.
$$}

\reponse{
Comme $G''-G=\delta$, on pose
$$
T=G\ast(\delta'-\delta'')=G'-G''.
$$
Calculons les dérivées de $G=H\sinh$ :
$$
G'=H\cosh,
$$
car $\sinh(0)=0$, puis
$$
G''=\delta+H\sinh,
$$
car $\cosh(0)=1$. Ainsi,
$$
T=H\cosh-\delta-H\sinh.
$$
Comme $\cosh t-\sinh t=e^{-t}$,
$$
\boxed{T=-\delta+H(t)e^{-t}}.
$$
Cette distribution est bien à support dans $\mathbb{R}_+$.
}

\item \question{Retrouver la solution précédente en utilisant la transformation de Fourier.}

\reponse{
La solution obtenue à la question précédente est une distribution tempérée. On peut donc appliquer la transformation de Fourier dans $\mathcal{S}'(\mathbb{R})$. En transformant l'équation, on obtient
$$
\bigl((i\xi)^2-1\bigr)\widehat T(\xi)
=
i\xi-(i\xi)^2.
$$
Ainsi,
$$
-(\xi^2+1)\widehat T(\xi)
=
i\xi+\xi^2,
$$
d'où
$$
\widehat T(\xi)
=-\frac{\xi^2+i\xi}{\xi^2+1}.
$$
On réécrit cette expression sous la forme
$$
\widehat T(\xi)
=-1+\frac{1-i\xi}{1+\xi^2}
=-1+\frac{1}{1+i\xi}.
$$
Or
$$
\widehat{\delta}=1
\qquad\text{et}\qquad
\widehat{H e^{-\cdot}}(\xi)=\frac{1}{1+i\xi}.
$$
Par inversion de Fourier,
$$
\boxed{T=-\delta+H(t)e^{-t}}.
$$
}

\end{enumerate}

}
